\subsection{Which classes account for the travel result?}
Averaging the class-specific travel distances in Appendix~\ref{app:cases} recovers the aggregate result in Table~\ref{tab:w1}. Relative to empirical transfer, diffusion reduces travel discrepancy by about 85.5 metres for buses, 20.1 metres for motorcycles, 73.1 metres for private cars and 18.7 metres for bicycles. For taxis it increases discrepancy by 79.0 metres. These are descriptive components of the prespecified macro average, not separate statistically significant class effects.

The finding for cars is supported by 33 reference conditions and that for buses by seven. The unfavorable taxi contribution receives the same weight despite its two references. It accounts for much of the difference between stating that the overall distribution improved and claiming improvement for every class. Giving every candidate equal weight across the entire test set would instead give private cars much greater influence. These would be different estimands. We use the stated macro measure consistently so the class-control question remains visible.

The stop fractions offer a second interpretation of the remaining discrepancy. For motorcycles, the observed set spends 24.6\% of sampled intervals below the stop threshold, while diffusion spends about 7.0\% and transfer about 7.4\%. For taxis the corresponding figures are 19.2\%, 2.2\% and 5.3\%. Both motion mechanisms underrepresent low-speed intervals for these particular references. The current model can generate zero increments, so their mere possibility under the decoder is insufficient to match their observed frequency. Whether this gap persists on a larger independent cohort remains open.
